% Options for packages loaded elsewhere
\PassOptionsToPackage{unicode}{hyperref}
\PassOptionsToPackage{hyphens}{url}
\PassOptionsToPackage{space}{xeCJK}
%
\documentclass[
  ignorenonframetext,
  aspectratio=169,
]{beamer}
\title{Dieckow 縦断16Sからの生態相互作用推定と検証}
\subtitle{組成時系列から符号付き相互作用行列 \(A\) へ ---
フィット・LOO-CV・誠実な検証}
\author{西岡佳祐 --- NIFE / SFB TRR-298}
\date{2026-06-03}

\usepackage{pgfpages}
\setbeamertemplate{caption}[numbered]
\setbeamertemplate{caption label separator}{: }
\setbeamercolor{caption name}{fg=normal text.fg}
\beamertemplatenavigationsymbolsempty
% Prevent slide breaks in the middle of a paragraph
\widowpenalties 1 10000
\raggedbottom
\setbeamertemplate{part page}{
  \centering
  \begin{beamercolorbox}[sep=16pt,center]{part title}
    \usebeamerfont{part title}\insertpart\par
  \end{beamercolorbox}
}
\setbeamertemplate{section page}{
  \centering
  \begin{beamercolorbox}[sep=12pt,center]{part title}
    \usebeamerfont{section title}\insertsection\par
  \end{beamercolorbox}
}
\setbeamertemplate{subsection page}{
  \centering
  \begin{beamercolorbox}[sep=8pt,center]{part title}
    \usebeamerfont{subsection title}\insertsubsection\par
  \end{beamercolorbox}
}
\AtBeginPart{
  \frame{\partpage}
}
\AtBeginSection{
  \ifbibliography
  \else
    \frame{\sectionpage}
  \fi
}
\AtBeginSubsection{
  \frame{\subsectionpage}
}
\usepackage{amsmath,amssymb}
\usepackage{lmodern}
\usepackage{iftex}
\ifPDFTeX
  \usepackage[T1]{fontenc}
  \usepackage[utf8]{inputenc}
  \usepackage{textcomp} % provide euro and other symbols
\else % if luatex or xetex
  \usepackage{unicode-math}
  \defaultfontfeatures{Scale=MatchLowercase}
  \defaultfontfeatures[\rmfamily]{Ligatures=TeX,Scale=1}
  \setmainfont[]{Noto Serif CJK JP}
  \setmonofont[]{Noto Sans Mono CJK JP}
  \ifXeTeX
    \usepackage{xeCJK}
    \setCJKmainfont[]{Noto Serif CJK JP}
  \fi
  \ifLuaTeX
    \usepackage[]{luatexja-fontspec}
    \setmainjfont[]{Noto Serif CJK JP}
  \fi
\fi
\usetheme[]{Madrid}
\usecolortheme{whale}
\usefonttheme{serif} % use mainfont rather than sansfont for slide text
% Use upquote if available, for straight quotes in verbatim environments
\IfFileExists{upquote.sty}{\usepackage{upquote}}{}
\IfFileExists{microtype.sty}{% use microtype if available
  \usepackage[]{microtype}
  \UseMicrotypeSet[protrusion]{basicmath} % disable protrusion for tt fonts
}{}
\makeatletter
\@ifundefined{KOMAClassName}{% if non-KOMA class
  \IfFileExists{parskip.sty}{%
    \usepackage{parskip}
  }{% else
    \setlength{\parindent}{0pt}
    \setlength{\parskip}{6pt plus 2pt minus 1pt}}
}{% if KOMA class
  \KOMAoptions{parskip=half}}
\makeatother
\usepackage{xcolor}
\IfFileExists{xurl.sty}{\usepackage{xurl}}{} % add URL line breaks if available
\IfFileExists{bookmark.sty}{\usepackage{bookmark}}{\usepackage{hyperref}}
\hypersetup{
  pdftitle={Dieckow 縦断16Sからの生態相互作用推定と検証},
  pdfauthor={西岡佳祐 --- NIFE / SFB TRR-298},
  hidelinks,
  pdfcreator={LaTeX via pandoc}}
\urlstyle{same} % disable monospaced font for URLs
\newif\ifbibliography
\usepackage{longtable,booktabs,array}
\usepackage{calc} % for calculating minipage widths
\usepackage{caption}
% Make caption package work with longtable
\makeatletter
\def\fnum@table{\tablename~\thetable}
\makeatother
\usepackage{graphicx}
\makeatletter
\def\maxwidth{\ifdim\Gin@nat@width>\linewidth\linewidth\else\Gin@nat@width\fi}
\def\maxheight{\ifdim\Gin@nat@height>\textheight\textheight\else\Gin@nat@height\fi}
\makeatother
% Scale images if necessary, so that they will not overflow the page
% margins by default, and it is still possible to overwrite the defaults
% using explicit options in \includegraphics[width, height, ...]{}
\setkeys{Gin}{width=\maxwidth,height=\maxheight,keepaspectratio}
% Set default figure placement to htbp
\makeatletter
\def\fps@figure{htbp}
\makeatother
\setlength{\emergencystretch}{3em} % prevent overfull lines
\providecommand{\tightlist}{%
  \setlength{\itemsep}{0pt}\setlength{\parskip}{0pt}}
\setcounter{secnumdepth}{-\maxdimen} % remove section numbering
\usepackage{amsmath,amssymb}
\newcommand{\sgn}{\operatorname{sgn}}
\newcommand{\relu}[1]{\left[#1\right]_{+}}
\ifLuaTeX
  \usepackage{selnolig}  % disable illegal ligatures
\fi

\begin{document}
\frame{\titlepage}

\begin{frame}{このデッキの位置づけ}
\protect\hypertarget{ux3053ux306eux30c7ux30c3ux30adux306eux4f4dux7f6eux3065ux3051}{}
本プロジェクト = 3 本柱 ＋ 空間拡張。データの流れ：

raw 16S → guild \(\varphi\) → gLV/Hamilton（＋符号 prior）→ LOO 検証 →
空間 PDE

デッキ構成（◀ = 本デッキ）：

\begin{itemize}
\tightlist
\item
  \textbf{Overview}（傘）--- 全体像と3本柱
\item
  \textbf{AGORA} --- 代謝 → 符号 prior（生態モデルの入力）
\item
  \textbf{Dieckow} --- in-vivo 相互作用の推定と検証（生態モデル本体） ◀
  \textbf{本デッキ}
\item
  \textbf{Network} --- 相互作用行列 \(A\) の構造解析
\item
  \textbf{Spatial-PDE} --- FISH 深さプロファイルの反応拡散
\item
  \textbf{FISH pipeline} --- .lif → 5 種深さ組成
\end{itemize}
\end{frame}

\begin{frame}{狙い}
\protect\hypertarget{ux72d9ux3044}{}
口腔バイオフィルムの\textbf{組成的 16S
時系列}から、\textbf{符号付き相互作用行列} \(A\) を
推定し、その推定をどこまで信用できるかを検証する。

\begin{itemize}
\tightlist
\item
  モデル：一般化 Lotka--Volterra (gLV) / Hamilton replicator（単体上）
\item
  制約：AGORA2 由来の \textbf{代謝符号
  prior}（別デッキ参照、ここでは簡潔に引用）
\item
  焦点：\textbf{フィット}・\textbf{LOO-CV}・モデル比較・\textbf{誠実な検証}
\end{itemize}

\vspace{0.4em}

\textbf{主張：} prior
は予測精度を買うのではなく、符号整合性・解釈可能性を与える。
データから独立に裏づくのは \textbf{cross-feeding 方向のみ}（後述）。
\end{frame}

\begin{frame}[fragile]{データと設計}
\protect\hypertarget{ux30c7ux30fcux30bfux3068ux8a2dux8a08}{}
Dieckow et al.~2024（ENA \textbf{PRJEB71108}）：\textbf{10 患者
\(\times\) 3 週}の縦断 16S。class レベルの \textbf{10
ギルド}に集約（\texttt{GUILD\_ORDER} を正典とする）。

\includegraphics[width=\textwidth,height=0.58\textheight]{dieckow_paper/figures/fig0_study_design.png}

ペリインプラント周囲の生着初期動態。各サンプルは組成
\(\varphi\in\Delta^{9}\) （\(\sum_i\varphi_i=1\)）として扱う。
\end{frame}

\begin{frame}{モデル --- gLV と replicator}
\protect\hypertarget{ux30e2ux30c7ux30eb-glv-ux3068-replicator}{}
絶対量 \(x_i\) の一般化 Lotka--Volterra：
\[\dot{x}_i = x_i\Big(b_i + \sum_{j=1}^{S} A_{ij}\,x_j\Big),\qquad S=10 .\]

単体上の組成 \(\varphi_i=x_i/\sum_k x_k\) に対する replicator / Hamilton
形：
\[\dot{\varphi}_i = \varphi_i\Big[(A\varphi+b)_i - \varphi^{\!\top}(A\varphi+b)\Big],
\qquad \sum_i\varphi_i=1 .\]

\begin{itemize}
\tightlist
\item
  \textbf{Hamilton} は \(A\) を\textbf{対称}に取る（線形ペイオフの
  replicator；Taylor \& Jonker 1978）。
\item
  \textbf{古典 gLV} は \(A\)
  が\textbf{非対称}。\(A_{ij}>0\)：促進、\(A_{ij}<0\)：抑制。
\end{itemize}
\end{frame}

\begin{frame}{推定は劣決定 \(\Rightarrow\) 代謝符号 prior}
\protect\hypertarget{ux63a8ux5b9aux306fux52a3ux6c7aux5b9a-rightarrow-ux4ee3ux8b1dux7b26ux53f7-prior}{}
パラメータ \(\mathcal{O}(S^2)\approx100\) に対し観測は
\(10\times3\)。存在量だけでは \(A\) の符号は識別不能 \(\Rightarrow\)
片側ヒンジで\textbf{符号違反のみ}を罰する：
\[\mathcal{L}(A,b)=\frac{1}{2\sigma^{2}}\sum_{t}\big\lVert\varphi^{\text{obs}}_{t}-\varphi^{\text{pred}}_{t}\big\rVert^{2}
\;+\; W\!\!\sum_{(i,j)\in\mathcal{M}}\!\!\left[-\,P_{ij}\,A_{ij}\right]_{+} .\]

\begin{itemize}
\tightlist
\item
  \(P_{ij}=\operatorname{sgn}(F_{ij})\) は AGORA cross-feeding
  フラックス由来（別デッキで導出）。
\item
  \(\left[-P_{ij}A_{ij}\right]_{+}\) は
  \(\operatorname{sgn}(A_{ij})=P_{ij}\) のとき \(0\) ---
  \textbf{大きさは不問}。
\item
  posterior は \textbf{TMCMC}（\(10^{4}\) 粒子）でサンプリング。
\end{itemize}
\end{frame}

\begin{frame}{フィットした相互作用行列 \(A\)}
\protect\hypertarget{ux30d5ux30a3ux30c3ux30c8ux3057ux305fux76f8ux4e92ux4f5cux7528ux884cux5217-a}{}
\includegraphics[width=\textwidth,height=0.6\textheight]{dieckow_paper/figures/fig1_A_matrix.png}

対角は\textbf{自己抑制} \(A_{ii}<0\)（資源制限／飽和）。 オフ対角に
Actinobacteria \(\leftrightarrow\) Bacilli の\textbf{促進}が現れる。
\end{frame}

\begin{frame}{軌道フィット（全 10 患者）}
\protect\hypertarget{ux8eccux9053ux30d5ux30a3ux30c3ux30c8ux5168-10-ux60a3ux8005}{}
\includegraphics[width=\textwidth,height=0.66\textheight]{dieckow_paper/figures/fig_fitting_results.png}

観測（点）と予測軌道（線）。患者ごとに \(b\) を、行列 \(A\)
は共有して推定。
\end{frame}

\begin{frame}{\(A_{ij}\) と代謝 \(F_{ij}\) の符号対応}
\protect\hypertarget{a_ij-ux3068ux4ee3ux8b1d-f_ij-ux306eux7b26ux53f7ux5bfeux5fdc}{}
\includegraphics[width=\textwidth,height=0.58\textheight]{dieckow_paper/figures/fig5b_aij_fij_scatter.png}

横軸 = 代謝フロー \(F_{ij}\)、縦軸 = 推定 \(\hat A_{ij}\)。 W=1.0
では制約セル \(\mathcal{M}\) 上で
\(\operatorname{sgn}(\hat A_{ij})=\operatorname{sgn}(F_{ij})\) が一致。
\end{frame}

\begin{frame}{LOO-CV の定義}
\protect\hypertarget{loo-cv-ux306eux5b9aux7fa9}{}
患者 \(p\) を抜いて行列を再推定し、抜いた患者の軌道で評価：
\[\mathrm{RMSE}=\sqrt{\frac{1}{NT}\sum_{t}\big\lVert\varphi^{\text{obs}}_{t}-\varphi^{\text{pred}}_{t}\big\rVert^{2}},\]
\[\mathrm{BC}=\frac{\sum_i\lvert\varphi^{\text{obs}}_i-\varphi^{\text{pred}}_i\rvert}{\sum_i\big(\varphi^{\text{obs}}_i+\varphi^{\text{pred}}_i\big)}
\quad(\text{Bray–Curtis 非類似度}).\]

10 患者すべてを順に hold-out し、平均する（leave-one-patient-out）。
\end{frame}

\begin{frame}{LOO モデル比較}
\protect\hypertarget{loo-ux30e2ux30c7ux30ebux6bd4ux8f03}{}
\includegraphics[width=\textwidth,height=0.44\textheight]{dieckow_paper/figures/fig2_loo_comparison.png}

\includegraphics[width=\textwidth,height=0.3\textheight]{dieckow_paper/figures/fig8_all_models_rmse_bc.png}

prior 層・モデル形・\(W\) を変えて LOO-RMSE / LOO-BC を比較。
\end{frame}

\begin{frame}{数値（要点）}
\protect\hypertarget{ux6570ux5024ux8981ux70b9}{}
ベストモデル \textbf{L1+L2+AGORA, \(W=1.0\)}：

\begin{longtable}[]{@{}ll@{}}
\toprule
指標 & 値 \\
\midrule
\endhead
train RMSE & \(0.0565\) \\
Pearson \(r\) & \(0.951\) \\
sign agreement & \(70/70\)（\(100\%\)） \\
\textbf{LOO-RMSE} & \(\mathbf{0.0504}\) \\
LOO-BC & \(0.1468\) \\
\bottomrule
\end{longtable}

\begin{itemize}
\tightlist
\item
  L1+L2 のみ：LOO-RMSE \(0.0516\)。
\item
  \textbf{prior-free gLV}：LOO-RMSE \(0.0455\)（より低い ---
  パラメータが多く代謝非制約）。
\item
  \textbf{MacArthur 量 prior は失敗}：SA \(4\text{–}8/70\)。
\end{itemize}

\vspace{0.2em}

\(\Rightarrow\) prior
の価値は予測精度ではなく\textbf{符号整合性・解釈可能性}。
\end{frame}

\begin{frame}{LOO での \(A\) 行列の安定性}
\protect\hypertarget{loo-ux3067ux306e-a-ux884cux5217ux306eux5b89ux5b9aux6027}{}
\includegraphics[width=\textwidth,height=0.6\textheight]{dieckow_paper/figures/fig7_loo_stability.png}

10 個の hold-out で再推定した行列間で、全ペアの\textbf{符号一致率
\(\geq 0.70\)}。 推定された符号構造は患者の入れ替えに対して頑健。
\end{frame}

\begin{frame}{批判的検証（誠実に）}
\protect\hypertarget{ux6279ux5224ux7684ux691cux8a3cux8aa0ux5b9fux306b}{}
prior を \textbf{OFF}（\(\alpha=0\)）にすると prior
は\textbf{全正}になり、naive な sign-agreement
は過大評価される。\textbf{ラベル permutation
検定}（\(n=10^{4}\)）で対照：

\begin{longtable}[]{@{}
  >{\raggedright\arraybackslash}p{(\columnwidth - 4\tabcolsep) * \real{0.33}}
  >{\raggedright\arraybackslash}p{(\columnwidth - 4\tabcolsep) * \real{0.33}}
  >{\raggedright\arraybackslash}p{(\columnwidth - 4\tabcolsep) * \real{0.33}}@{}}
\toprule
\begin{minipage}[b]{\linewidth}\raggedright
モデル
\end{minipage} & \begin{minipage}[b]{\linewidth}\raggedright
cross-feeding 方向
\end{minipage} & \begin{minipage}[b]{\linewidth}\raggedright
競争方向
\end{minipage} \\
\midrule
\endhead
\textbf{Hamilton \(\alpha=0\)} & \(\mathbf{78.6\%}\) (11/14) vs random
\(37.7\%\), \(p=4\times10^{-4}\), \(z=+3.79\) & \(\approx\) chance \\
gLV & \(41\%\)（null） & null \\
\bottomrule
\end{longtable}

\includegraphics[width=\textwidth,height=0.3\textheight]{results/dieckow_cr/loo_alpha_comparison_hamilton_noagora_glv_noagora.png}

\textbf{検証されるのは cross-feeding 方向のみ。競争方向は支持されない。}
\end{frame}

\begin{frame}{2 コホート再現}
\protect\hypertarget{ux30b3ux30dbux30fcux30c8ux518dux73fe}{}
Botelho 2021（\textbf{PRJNA725874}、15 患者 \(\times\) 7 時点）を
\textbf{prior 抜き}で 独立にフィットし、Dieckow と強ペアの符号を比較：

\includegraphics[width=\textwidth,height=0.46\textheight]{results/botelho_validation/fig_botelho_A_comparison_noprior.png}

強ペアの有向符号が \textbf{\(89\%\)}（\(8/9\)
上三角）で一致（\(p\approx0.02\)）。 特に \textbf{Actinobacteria
軸}が両コホートで整合。
\end{frame}

\begin{frame}{誠実な解釈}
\protect\hypertarget{ux8aa0ux5b9fux306aux89e3ux91c8}{}
\begin{itemize}
\tightlist
\item
  \textbf{AGORA prior 自体は 16S 力学で再現されない} \(\Rightarrow\)
  それは \textbf{モデリング選択}であり、データ確認済みの事実ではない。
\item
  データから独立に裏づくのは：

  \begin{enumerate}
  \tightlist
  \item
    \textbf{cross-feeding 方向}（Hamilton, \(p=4\times10^{-4}\)）、
  \item
    \textbf{2 コホート強ペア}符号（Dieckow \(\times\) Botelho,
    \(89\%\)）。
  \end{enumerate}
\item
  競争方向・prior の大きさ・全体行列は検証されていない。
\end{itemize}

\vspace{0.3em}

過大に語らない：符号の\textbf{一部}が支持され、残りは仮説に留まる。
\end{frame}

\begin{frame}{外部臨床検証（進行中）}
\protect\hypertarget{ux5916ux90e8ux81e8ux5e8aux691cux8a3cux9032ux884cux4e2d}{}
Guild Dysbiosis Index を Joshi 2025
コホート（\textbf{PRJNA1192962}）と対照する計画：

\includegraphics[width=\textwidth,height=0.46\textheight]{dieckow_paper/figures/fig3_joshi_attractor.png}

\textcolor{red}{\textbf{予備的}：臨床メタデータ待ち。本図は暫定であり、健常／疾患
ラベルとの突き合わせは未完了。}
\end{frame}

\begin{frame}{結論}
\protect\hypertarget{ux7d50ux8ad6}{}
\begin{enumerate}
\tightlist
\item
  組成 16S 時系列 \(\to\) 符号付き \(A\) を gLV / Hamilton replicator
  で推定。
\item
  劣決定性は \textbf{AGORA 代謝符号
  prior}（片側ヒンジ）で緩和；posterior は TMCMC。
\item
  ベスト：L1+L2+AGORA \(W=1.0\) で \textbf{LOO-RMSE \(0.0504\)}、SA
  \(100\%\)。 ただし prior-free gLV は \(0.0455\) --- prior
  は精度でなく\textbf{解釈可能性}を買う。
\item
  \textbf{LOO で符号は安定}（\(\geq0.70\)）；MacArthur 量 prior は失敗。
\item
  誠実な検証：\textbf{cross-feeding 方向}（\(p=4\times10^{-4}\)）と
  \textbf{2
  コホート強ペア}（\(89\%\)）のみが独立に支持。臨床検証は進行中。
\end{enumerate}
\end{frame}

\end{document}
